% ============================================================
% PARTE VI - GLI ACCORDI DI SETTIMA
% CAPITOLO 22 - LA SETTIMA DI DOMINANTE
% ============================================================
\part{Gli accordi di settima}

\chapter{La settima di dominante}
\label{chap:settima_dominante}

Aggiungendo una quarta nota alla triade di dominante si ottiene la \textbf{settima di dominante} (\sigla{V}{7}{}), un accordo di forte tensione che spinge verso la risoluzione.

\section{Genesi e struttura}
\label{sec:genesi_v7}

L'accordo si costruisce aggiungendo una terza minore sopra la quinta della triade di dominante. In Do maggiore la dominante è Sol-Si-Re; una terza minore sopra il Re dà il Fa, e l'accordo completo è Sol-Si-Re-Fa.

\begin{definizione}[Il tritono]
	\label{def:tritono}
	Il \textbf{tritono} è un intervallo di tre toni interi, che può presentarsi come quarta aumentata o come quinta diminuita. È l'intervallo più dissonante del sistema diatonico ed era chiamato storicamente \textit{diabolus in musica}.
\end{definizione}

La \sigla{V}{7}{} contiene un tritono tra la sua terza (la sensibile) e la sua settima: in Do maggiore tra il \textbf{Si} e il \textbf{Fa}. Questa dissonanza interna è la fonte dell'instabilità dell'accordo e del suo bisogno di risoluzione. La base fisica della tensione del tritono, legata ai battimenti, è nell'Approfondimento \ref{phys:battimenti_tritono}.

\begin{esempio}[La \sigla{V}{7}{} e il suo tritono in Do maggiore]
	\label{ex:tritono_v7}
	\musica{es_cap09_01}
\end{esempio}

\begin{regola}[Preparazione e risoluzione della dissonanza]
	\label{reg:preparazione_risoluzione}
	Ogni nota dissonante segue il ciclo di preparazione, percussione e risoluzione (Capitolo \ref{chap:teoria_dissonanza}). Per la settima di dominante la preparazione viene spesso omessa, data la familiarità dell'accordo, ma la \textbf{risoluzione è obbligatoria}.
\end{regola}

\section{Risoluzione naturale \sigla{V}{7}{} - I}
\label{sec:risoluzione_v7}

La risoluzione naturale della settima di dominante è sulla tonica. Le due note del tritono hanno movimenti obbligati.

\begin{regola}[Vettori obbligati della \sigla{V}{7}{}]
	\label{reg:vettori_obbligati}
	Nella risoluzione \sigla{V}{7}{} $\rightarrow$ I:
	\begin{itemize}
		\item la \textbf{terza} dell'accordo (la sensibile) \textbf{sale} di semitono sulla tonica (in Do: Si $\nearrow$ Do);
		\item la \textbf{settima} \textbf{scende} di grado sulla terza della tonica (in Do: Fa $\searrow$ Mi).
	\end{itemize}
\end{regola}

Quattro note devono risolvere su una triade che ne ha tre: la soluzione sta nelle note omesse.

\begin{definizione}[\sigla{V}{7}{} completa e incompleta]
	\label{def:v7_completa_incompleta}
	Ci sono due modi di risolvere la \sigla{V}{7}{}:
	\begin{enumerate}
		\item \textbf{\sigla{V}{7}{} completa $\rightarrow$ I incompleta:} si usano tutte e quattro le note. Seguendo i vettori obbligati, la tonica risultante è senza quinta e ha la fondamentale tre volte.
		\item \textbf{\sigla{V}{7}{} incompleta $\rightarrow$ I completa:} si omette la quinta della \sigla{V}{7}{} (la nota meno caratteristica) e si raddoppia la fondamentale. La tonica è completa, con fondamentale, terza e quinta.
	\end{enumerate}
\end{definizione}

\begin{esempio}[Le due risoluzioni della \sigla{V}{7}{}]
	\label{ex:due_risoluzioni_v7}
	\textbf{A. \sigla{V}{7}{} completa $\rightarrow$ I incompleta}
	\musica{es_cap09_02}

	\vspace{1cm}
	\textbf{B. \sigla{V}{7}{} incompleta $\rightarrow$ I completa}
	\musica{es_cap09_03}
\end{esempio}

\section{\sigla{V}{7}{} in stato fondamentale (cifra 7)}
\label{subsec:v7_fondamentale}

La settima di dominante si presenta in quattro stati, uno fondamentale e tre rivolti. In ogni stato i vettori del tritono (sensibile $\nearrow$ tonica, settima $\searrow$ terza) restano il principio guida. Gli stati si distinguono con la numerica sintetica (Sezione \ref{subsec:cifratura_accordi}): 7 per lo stato fondamentale, 6/5, 4/3 e 2 per i rivolti. Negli esempi ogni stato è mostrato in Do maggiore e in La minore (scala armonica). Le note comuni a due accordi sono scritte come due semiminime \textbf{legate}; le frecce colorate indicano il moto delle voci che cambiano nota (blu soprano, verde contralto, arancione tenore, rosso basso).

\begin{itemize}
    \item \textbf{Al basso:} la fondamentale (Sol in Do maggiore).
    \item \textbf{Risoluzione:} $\rightarrow$ I in stato fondamentale.
    \item \textbf{Condotta:}
    \begin{itemize}
        \item Basso (Sol): sale di quarta a Do.
        \item Sensibile (Si): sale a Do.
        \item Settima (Fa): scende a Mi.
        \item Quinta (Re): scende a Do (\sigla{V}{7}{} completa $\rightarrow$ I incompleta). Se la quinta è omessa, la fondamentale raddoppiata resta ferma e diventa la quinta di una tonica completa.
    \end{itemize}
\end{itemize}
Sulla completezza si veda la Definizione \ref{def:v7_completa_incompleta}.

\begin{esempio}[\sigla{V}{7}{} in stato fondamentale in La minore]
	\label{ex:v7_fond_minore}
	La sensibile Sol$\sharp$ sale a La, la settima Re scende a Do. A sinistra \sigla{V}{7}{} completa $\rightarrow$ i incompleta, a destra \sigla{V}{7}{} incompleta $\rightarrow$ i completa.
	\musica{es_cap09_05}
\end{esempio}

\section{\sigla{V}{6}{5}, primo rivolto (cifra 6/5)}
\label{subsec:v65}
\begin{itemize}
    \item \textbf{Al basso:} la terza dell'accordo (Si in Do maggiore), cioè la \textbf{sensibile}.
    \item \textbf{Risoluzione:} $\rightarrow$ I \textbf{in stato fondamentale}, con il basso che sale di semitono dalla sensibile alla tonica.
    \item \textbf{Condotta:}
    \begin{itemize}
        \item Basso (Si): sale obbligatoriamente a Do.
        \item Settima (Fa): scende obbligatoriamente a Mi.
        \item Quinta (Re): scende a Do; la fondamentale (Sol) resta ferma e diventa la quinta della tonica.
    \end{itemize}
\end{itemize}

\begin{esempio}[\sigla{V}{6}{5} $\rightarrow$ I]
	\label{ex:v65}
	\musica{es_cap09_06}
\end{esempio}

\section{\sigla{V}{4}{3}, secondo rivolto (cifra 4/3)}
\label{subsec:v43}
\begin{itemize}
    \item \textbf{Al basso:} la quinta dell'accordo (Re in Do maggiore).
    \item \textbf{Risoluzione:} $\rightarrow$ I (il basso scende a Do) oppure $\rightarrow$ \sigla{I}{6}{} (il basso sale a Mi).
    \item \textbf{Condotta:}
    \begin{itemize}
        \item Basso (Re): scende o sale di grado.
        \item Sensibile (Si): sale a Do.
        \item Settima (Fa): scende a Mi.
        \item Fondamentale (Sol): resta ferma.
    \end{itemize}
    \item \textbf{Eccezione:} nella risoluzione su \sigla{I}{6}{} la settima può \textbf{salire} a Sol, in decime parallele con il basso (Re-Mi): così si evita di raddoppiare il Mi, terza della tonica, al basso e in una voce superiore. È l'unico caso in cui la settima di dominante non scende.
\end{itemize}

\begin{esempio}[\sigla{V}{4}{3} $\rightarrow$ I]
	\label{ex:v43}
	\musica{es_cap09_07}
\end{esempio}

\section{\sigla{V}{2}{}, terzo rivolto (cifra 2 o 4/2)}
\label{subsec:v2}
\begin{itemize}
    \item \textbf{Al basso:} la settima dell'accordo (Fa in Do maggiore).
    \item \textbf{Risoluzione:} $\rightarrow$ \sigla{I}{6}{}, con il basso che scende a Mi. È l'unica risoluzione possibile, perché la settima al basso deve scendere di grado.
    \item \textbf{Condotta:}
    \begin{itemize}
        \item Basso (Fa): scende obbligatoriamente a Mi.
        \item Sensibile (Si): sale a Do.
        \item Fondamentale (Sol): resta ferma; la tonica in primo rivolto ha il Do raddoppiato.
    \end{itemize}
\end{itemize}

\begin{esempio}[\sigla{V}{2}{} $\rightarrow$ \sigla{I}{6}{}]
	\label{ex:v2}
	\musica{es_cap09_08}
\end{esempio}

Nel modo minore la \sigla{V}{7}{} ha un'intensità particolare, perché si sommano la tensione della sensibile alzata e quella della settima (Approfondimento \ref{phys:doppia_tensione_v7}). Una tabella riassuntiva dei quattro stati è nel Capitolo \ref{chap:riepilogo_vi}.
